\documentclass[11pt]{article}

\usepackage[letterpaper,margin=1in]{geometry}
\usepackage[T1]{fontenc}
\usepackage[utf8]{inputenc}
\usepackage{lmodern}
\usepackage{microtype}
\usepackage{amsmath,amssymb,mathtools,bm}
\usepackage{slashed}
\usepackage{enumitem}
\usepackage{xcolor}
\usepackage{booktabs}
\usepackage{array}
\usepackage{tocloft}
\usepackage{hyperref}
\usepackage[most]{tcolorbox}
\usepackage{fancyhdr}
\usepackage{titlesec}
\usepackage{parskip}
\usepackage{needspace}
\setlength{\cftsecnumwidth}{2.4em}

\definecolor{qftblue}{HTML}{245A8D}
\definecolor{qftgreen}{HTML}{2F6B4F}
\definecolor{qftred}{HTML}{A23B3B}
\definecolor{qftgray}{HTML}{F4F6F8}

\hypersetup{
  colorlinks=true,
  linkcolor=qftblue,
  urlcolor=qftblue,
  citecolor=qftgreen,
  pdftitle={PHYS 662 Lecture 15: Canonical Quantization of Free Fields},
  pdfauthor={Nima Lashkari}
}

\pagestyle{fancy}
\fancyhf{}
\lhead{PHYS 662: Quantum Field Theory I}
\rhead{Lecture 15}
\cfoot{\thepage}
\setlength{\headheight}{14pt}

\titleformat{\section}{\Large\bfseries\color{qftblue}}{\thesection.}{0.6em}{}
\titleformat{\subsection}{\large\bfseries\color{qftgreen}}{\thesubsection}{0.6em}{}
\titleformat{\subsubsection}{\normalsize\bfseries\color{qftred}}{\thesubsubsection}{0.6em}{}

\setlist[itemize]{leftmargin=2em,itemsep=0.25em,topsep=0.4em}
\setlist[enumerate]{leftmargin=2em,itemsep=0.25em,topsep=0.4em}

\newcommand{\dd}{\mathrm{d}}
\newcommand{\R}{\mathbb{R}}
\newcommand{\Lag}{\mathcal{L}}
\newcommand{\cM}{\mathcal{M}}
\newcommand{\SO}{\mathrm{SO}}
\newcommand{\OO}{\mathrm{O}}

\newtcolorbox{keyidea}[1][]{
  enhanced, breakable, colback=qftgray, colframe=qftblue,
  boxrule=0.8pt, arc=1.5mm, title={#1}, fonttitle=\bfseries,
  coltitle=white
}
\newtcolorbox{remarkbox}[1][]{
  enhanced, breakable, colback=white, colframe=qftgreen,
  boxrule=0.7pt, arc=1.5mm, title={#1}, fonttitle=\bfseries,
  coltitle=white
}
\newtcolorbox{warningbox}[1][]{
  enhanced, breakable, colback=white, colframe=qftred,
  boxrule=0.7pt, arc=1.5mm, title={#1}, fonttitle=\bfseries,
  coltitle=white
}


\newcommand{\Th}{\Theta}
\newcommand{\Eul}{\mathcal E}
\newcommand{\Lie}{\mathscr L}

\newcommand{\C}{\mathbb C}
\newcommand{\Z}{\mathbb Z}
\newcommand{\SU}{\mathrm{SU}}
\newcommand{\UU}{\mathrm{U}}
\newcommand{\GL}{\mathrm{GL}}
\newcommand{\SL}{\mathrm{SL}}
\newcommand{\Sp}{\mathrm{Sp}}
\newcommand{\tr}{\operatorname{tr}}
\newcommand{\Ad}{\operatorname{Ad}}
\newcommand{\ad}{\operatorname{ad}}
\newcommand{\lieg}{\mathfrak g}
\newcommand{\lieh}{\mathfrak h}
\newcommand{\su}{\mathfrak{su}}
\newcommand{\so}{\mathfrak{so}}
\newcommand{\slie}{\mathfrak{sl}}


\newcommand{\cH}{\mathcal H}
\newcommand{\Spin}{\operatorname{Spin}}
\newcommand{\id}{\mathrm{id}}
\newcommand{\Lor}{\SO^+(1,3)}
\newcommand{\eps}{\varepsilon}
\newcommand{\one}{\mathbf 1}

\newcommand{\smu}{s}
\newcommand{\tmu}{\widetilde s}
\newcommand{\Eps}{E}
\newcommand{\pslash}{\slashed p}
\newcommand{\ISO}{\mathrm{ISO}}
\newcommand{\pL}{P_L}
\newcommand{\pR}{P_R}
\newcommand{\hc}{\mathrm{h.c.}}
\newcommand{\diag}{\operatorname{diag}}
\newcommand{\vol}{\mathrm{vol}}
\newcommand{\supp}{\operatorname{supp}}

\newcommand{\ddPi}{\dd\Pi}
\newcommand{\vac}{\lvert0\rangle}
\newcommand{\vbra}{\langle0\rvert}
\newcommand{\ii}{\mathrm i}
\newcommand{\Dc}{\mathcal C}
\newcommand{\Dp}{\mathcal P}
\newcommand{\Dt}{\mathcal T}
\begin{document}
\pagenumbering{gobble}
\begin{titlepage}
\centering
\vspace*{1.2in}
{\scshape\Large PHYS 662 --- Quantum Field Theory I\par}
\vspace{0.35in}
{\Huge\bfseries\color{qftblue} Lecture 15\par}
\vspace{0.18in}
{\LARGE\bfseries Canonical Quantization of Free Fields\par}
\vspace{0.16in}
{\large\color{qftgreen} Scalar and Dirac Fields, Fock Space, and Antiparticles\par}
\vspace{0.6in}
{\large Nima Lashkari\par}
{\large Fall 2026\par}
\vfill
\begin{minipage}{0.86\textwidth}\small
We construct quantum scalar and Dirac fields from their classical modes. The comparison makes the roles of normalization, statistics, vacuum choice, energy, and conserved charge explicit.

\medskip
Course page: \url{https://nimalashkari.com/quantum-field-theory-2026-part-i/}
\end{minipage}
\vspace{0.5in}
\end{titlepage}
\pagenumbering{roman}
\setcounter{tocdepth}{1}\begingroup\small\tableofcontents\endgroup
\newpage
\pagenumbering{arabic}

\section*{Roadmap}
\addcontentsline{toc}{section}{Roadmap}
Lectures 11-14 supplied the particle representations and classical wavefunctions.
We now quantize both scalar and Dirac fields. The common steps are to fix
canonical data, normalize the modes, choose a vacuum, and derive energy and
charge. Their different oscillator algebras produce Bose symmetry and Fermi
exclusion. Lecture 16 will use these constructions to study discrete
symmetries, locality, and propagators.

\begin{keyidea}[One normalization throughout]
Use \(\eta=\diag(1,-1,-1,-1)\), \(\hbar=c=1\),
\(p\cdot x=E_{\bm p}t-\bm p\cdot\bm x\), and
\[
 E_{\bm p}=\sqrt{\bm p^2+m^2},\qquad
 \dd\Pi_p=\frac{\dd^d p}{(2\pi)^d2E_{\bm p}}.
\]
There are \(d\) spatial dimensions; set \(d=3\) unless indicated otherwise.
The covariantly normalized operators obey
\([a(\bm p),a^\dagger(\bm q)]=(2\pi)^d2E_{\bm p}
\delta^{(d)}(\bm p-\bm q)\). The measure and this commutator must be
changed together if another normalization is used.
\end{keyidea}

\Needspace{7\baselineskip}
\section{From a lattice of oscillators to a field}

\subsection{What survives the continuum limit?}
The harmonic chain gives a regulated example of the construction. Let
\(q_j\) be the displacement of an atom of mass \(M\), with lattice
spacing \(\ell\), periodic circumference \(L=N\ell\), nearest-neighbor
spring constant \(\kappa\), and an on-site restoring constant \(\beta\):
\begin{equation}
 L_{\rm chain}=\sum_{j=0}^{N-1}\left[
 \frac M2\dot q_j^2-\frac\kappa2(q_{j+1}-q_j)^2-\frac\beta2q_j^2\right].
 \label{s-eq:chain}
\end{equation}
A discrete Fourier transform diagonalizes this quadratic form. Substituting
\(q_j\propto e^{i(kj\ell-\omega t)}\) into the equations of motion gives
\begin{equation}
 \omega_k^2=\frac\beta M+\frac{4\kappa}{M}\sin^2\frac{k\ell}{2}
 =\Omega_0^2+c_s^2k^2+O(\ell^2k^4),\qquad
 c_s=\ell\sqrt{\frac\kappa M},\quad \Omega_0^2=\frac\beta M.
 \label{s-eq:dispersion}
\end{equation}
Only long wavelengths, \(|k|\ell\ll1\), have the relativistic form.
With \(\phi=\sqrt{M/\ell}\,q\), the continuum Lagrangian is
\[
 L=\frac12\int\dd x\,[\dot\phi^2-c_s^2(\partial_x\phi)^2-\Omega_0^2\phi^2].
\]
Choose time and length units with \(c_s=1\), and write the gap as \(m\).
A continuum limit holds these continuum parameters fixed while the
microscopic parameters scale with \(\ell\). Merely sending \(\ell\) to
zero at fixed \(M,\kappa\) would also send the wave speed to zero.

\subsection{Canonical data for the Klein--Gordon field}

The relativistic theory in \(d+1\) dimensions is
\begin{equation}
 \Lag=\frac12\partial_\mu\phi\partial^\mu\phi-\frac12m^2\phi^2,
 \quad \pi=\frac{\partial\Lag}{\partial \dot\phi}=\dot\phi,\quad
 H=\frac12\int\dd^d x\,[\pi^2+(\bm\nabla\phi)^2+m^2\phi^2].
 \label{s-eq:realaction}
\end{equation}
The field equation is \((\Box+m^2)\phi=0\).
Canonical quantization promotes the field and its momentum to operators:
\begin{equation}
 [\phi(t,\bm x),\pi(t,\bm y)]=i\delta^{(d)}(\bm x-\bm y),\qquad
 [\phi(t,\bm x),\phi(t,\bm y)]=[\pi(t,\bm x),\pi(t,\bm y)]=0.
 \label{s-eq:CCR}
\end{equation}
These are initial-data relations at equal time. They do not say that
fields commute at arbitrary spacetime separation. To simplify notation, hats will be omitted.

\Needspace{7\baselineskip}
\section{Mode expansion and oscillator algebra}
\subsection{Deriving the normalization}

At \(t=0\), use spatial Fourier transforms with
\(\widetilde\phi(\bm p)=\int\dd^d x\,e^{-i\bm p\cdot\bm x}\phi(\bm x)\),
and similarly for \(\pi\). Reality implies
\(\widetilde\phi(\bm p)^\dagger=\widetilde\phi(-\bm p)\), and the canonical
commutator becomes
\[
 [\widetilde\phi(\bm p),\widetilde\pi(\bm q)]
   =i(2\pi)^d\delta^{(d)}(\bm p+\bm q).
\]
Define
\begin{equation}
 a(\bm p)=E_{\bm p}\widetilde\phi(\bm p)+i\widetilde\pi(\bm p),\qquad
 \widetilde\phi(\bm p)=\frac{a(\bm p)+a^\dagger(-\bm p)}{2E_{\bm p}},
 \quad \widetilde\pi(\bm p)=\frac{a(\bm p)-a^\dagger(-\bm p)}{2i}.
 \label{s-eq:inverse}
\end{equation}
The two cross terms in \([a,a^\dagger]\) each contribute
\(E_{\bm p}(2\pi)^d\delta^{(d)}(\bm p-\bm q)\). Thus
\begin{equation}
 [a(\bm p),a^\dagger(\bm q)]=(2\pi)^d2E_{\bm p}\delta^{(d)}(\bm p-\bm q),
 \qquad [a,a]=[a^\dagger,a^\dagger]=0.
 \label{s-eq:modeCCR}
\end{equation}
Solving the field equation with these initial data yields
\begin{align}
 \phi(x)&=\int\dd\Pi_p\,[a(\bm p)e^{-ip\cdot x}+a^\dagger(\bm p)e^{ip\cdot x}],
 \label{s-eq:realmode}\\
 \pi(x)&=-i\int\dd\Pi_p\,E_{\bm p}
 [a(\bm p)e^{-ip\cdot x}-a^\dagger(\bm p)e^{ip\cdot x}].
 \label{s-eq:pimode}
\end{align}
Checking \eqref{s-eq:CCR} now gives two identical Fourier integrals, each
with coefficient \(i/2\). Their sum is \(i\delta^{(d)}\), as required.

\begin{remarkbox}[The other common convention]
Writing \(c(\bm p)=a(\bm p)/\sqrt{2E_{\bm p}}\) gives
\([c(\bm p),c^\dagger(\bm q)]=(2\pi)^d\delta^{(d)}(\bm p-\bm q)\).
Then the mode measure is
\(\dd^d p/[(2\pi)^d\sqrt{2E_{\bm p}}]\).
Both conventions describe the same field. A factor \(1/(2E)\) in the
expansion together with the \(c\)-commutator would not reproduce the
canonical relation. The advantage of the notation $a(\bm p)$ is its Lorentz covariance, as we will see $U(\Lambda)a^\dagger(\bm p) U(\Lambda)^{-1}=a^\dagger(\Lambda \bm p)$. The advantage of the notation $c(\bm p)$ is that it is normalized similar to simple harmonic oscillators, so it is more convenient to use in constructing the occupation number basis in finite volume that breaks Lorentz invariance, as we see below.
\end{remarkbox}

\subsection{The Hamiltonian and the sign of the frequency}
Inserting the modes into \eqref{s-eq:realaction}, spatial integration sets
opposite momenta in the \(aa\) and \(a^\dagger a^\dagger\) terms.
Their coefficient is \(-E_{\bm p}^2+\bm p^2+m^2=0\). The remaining terms give
\begin{equation}
 H=\frac12\int\dd\Pi_p\,E_{\bm p}
 [a^\dagger(\bm p)a(\bm p)+a(\bm p)a^\dagger(\bm p)].
 \label{s-eq:Hsym}
\end{equation}
Consequently \([H,a^\dagger(\bm p)]=E_{\bm p}a^\dagger(\bm p)\) and
\([H,a(\bm p)]=-E_{\bm p}a(\bm p)\). The creation operator multiplies
\(e^{+ip\cdot x}\), yet creates a positive-energy excitation. The sign
of a classical Fourier frequency is not, by itself, the energy of the
state created by its coefficient.

\Needspace{7\baselineskip}
\section{Vacuum, Fock space, and relativistic states}
\subsection{Regulate before counting occupation numbers}
In a periodic box of volume \(V\), momenta are discrete and
\[
 \int\frac{\dd^d p}{(2\pi)^d}\longrightarrow\frac1V\sum_{\bm p},
 \qquad (2\pi)^d\delta^{(d)}(\bm p-\bm q)\longrightarrow V\delta_{\bm p\bm q}.
\]
Then \(c_{\bm p}=a(\bm p)/\sqrt{2E_{\bm p}V}\) satisfies
\([c_{\bm p},c_{\bm q}^\dagger]=\delta_{\bm p\bm q}\). Choose a normalized
vacuum annihilated by every \(c_{\bm p}\). The occupation basis is
\begin{equation}
 |\{n_{\bm p}\}\rangle=\prod_{\bm p}
 \frac{(c_{\bm p}^\dagger)^{n_{\bm p}}}{\sqrt{n_{\bm p}!}}\vac,
 \qquad n_{\bm p}=0,1,2,\ldots.
 \label{s-eq:fock}
\end{equation}
Creation operators commute, so exchange of two identical excitations
leaves their state unchanged. In the continuum the construction is
\(\mathcal F_B(\cH_1)=\bigoplus_{n=0}^{\infty}\operatorname{Sym}^n\cH_1\).
The number operator is
\begin{equation}
 N=\int\dd\Pi_p\,a^\dagger(\bm p)a(\bm p).
 \label{s-eq:N}
\end{equation}
An expression such as \(a^\dagger(\bm p)a(\bm p)\) at a fixed continuum
momentum is a density, not a single normalized oscillator number.
The box or a normalized wave packet makes that distinction precise.
For a massless scalar in finite volume, the exactly zero momentum mode
is a free-particle degree of freedom, not an oscillator with a normalizable
Gaussian ground state. Use \(m>0\), or treat that mode separately.

\subsection{The one-particle sector of Lecture 14}
Set \(|\bm p\rangle=a^\dagger(\bm p)\vac\). Then
\begin{equation}
 \langle\bm q|\bm p\rangle=(2\pi)^d2E_{\bm p}\delta^{(d)}(\bm p-\bm q),
 \qquad I_1=\int\dd\Pi_p\,|\bm p\rangle\langle\bm p|.
 \label{s-eq:1particle}
\end{equation}
The invariant measure follows by doing the \(p^0\) integral in
\[
 \frac{\dd^{d+1}p}{(2\pi)^d}\,\theta(p^0)\delta(p^2-m^2).
\]
The derivative of \((p^0)^2-E_{\bm p}^2\) supplies the factor \(2E_{\bm p}\).
For a spin-zero particle,
\(U(\Lambda)|\bm p\rangle=|\boldsymbol{\Lambda p}\rangle\).
An invariant vacuum then gives
\(U(\Lambda)a^\dagger(\bm p)U(\Lambda)^{-1}=a^\dagger(\boldsymbol{\Lambda p})\)
and \(U(\Lambda)\phi(x)U(\Lambda)^{-1}=\phi(\Lambda x)\).
The equivalent form with \(U^{-1}\) on the left has \(\Lambda^{-1}x\).

For \(|f\rangle=\int\dd\Pi_p\,f(\bm p)|\bm p\rangle\),
\(\|f\|^2=\int\dd\Pi_p|f(\bm p)|^2\). The field connects this state to a
positive-frequency wavefunction:
\begin{equation}
 \vbra\phi(x)|\bm p\rangle=e^{-ip\cdot x},\qquad
 \vbra\phi(x)|f\rangle=\int\dd\Pi_p f(\bm p)e^{-ip\cdot x}.
 \label{s-eq:wavefunction}
\end{equation}
The formal state \(\phi(x)\vac\) is not a normalized position eigenstate.
It must be smeared, and its overlap with \(\phi(y)\vac\) is a two-point
function rather than a spatial delta function. Lecture 16 explains its
nonzero spacelike overlap.

\Needspace{7\baselineskip}
\section{Vacuum energy, normal ordering, and ultraviolet scales}
\subsection{Volume divergence versus short-distance divergence}
With a finite box and a momentum cutoff, \eqref{s-eq:Hsym} reads
\begin{equation}
 H=\sum_{\bm p}E_{\bm p}\left(c_{\bm p}^\dagger c_{\bm p}+\frac12\right),
 \qquad E_0=\frac12\sum_{\bm p}E_{\bm p}.
 \label{s-eq:vacuumenergy}
\end{equation}
Taking large volume gives \(E_0=V\rho_0\), where in three spatial dimensions
\begin{align}
 \rho_0(\Lambda)&=\frac12\int_{|\bm p|<\Lambda}
   \frac{\dd^3p}{(2\pi)^3}\sqrt{\bm p^2+m^2}
   =\frac{1}{4\pi^2}\int_0^\Lambda\dd p\,p^2\sqrt{p^2+m^2},\notag\\
 &=\frac{\Lambda^4}{16\pi^2}
   +\frac{m^2\Lambda^2}{16\pi^2}
   -\frac{m^4}{32\pi^2}\log\frac{2\Lambda}{m}+O(m^4).
 \label{s-eq:rho}
\end{align}
The factor \(V\) is extensivity in infinite volume. The divergence of
\(\rho_0\) as \(\Lambda\to\infty\) is ultraviolet. These are different
limits; the volume factor is not a divergence of the small-momentum
integral for this massive theory. In \(d\) dimensions a radial ultraviolet
integral behaves as \(\int^\Lambda\dd p\,p^{d-1}f(p)\), so \(f(p)\sim p^{-d}\)
gives a logarithm.

Normal ordering relative to this vacuum places every creation operator
to the left without adding the commutators generated by the reordering.
For the free Hamiltonian it defines
\begin{equation}
 H_N\equiv:H:=H-E_0=\int\dd\Pi_p E_{\bm p}a^\dagger(\bm p)a(\bm p),
 \qquad \bm P=\int\dd\Pi_p\bm p\,a^\dagger(\bm p)a(\bm p).
 \label{s-eq:normalH}
\end{equation}
The shift leaves energy differences and Heisenberg evolution unchanged.
It is a choice of vacuum-energy reference in flat spacetime; coupling to
gravity requires a separate treatment of the vacuum-energy parameter.

\subsection{What cutoff independence means}
Fields at points are operator-valued distributions. Products such as
\(\phi(x)^2\) require a prescription; free-field normal ordering supplies
one, for example \(:\phi^2:=\phi^2-\vbra\phi^2\vac\) with a common regulator.
Normal ordering alone does not renormalize general interacting theories.

\begin{warningbox}[A regulator is not the same as a physical high-energy scale]
Renormalized predictions become independent of an arbitrary regulator
only after parameters are adjusted to fixed physical conditions.
In an effective field theory, short-distance physics enters masses,
couplings, and coefficients of higher-dimension operators. Once those
coefficients are matched, truncation errors are suppressed by powers of
energy divided by the physical high-energy scale, under the assumptions
of the effective expansion. Large cutoff by itself does not remove this
dependence. A sharp spatial cutoff also selects a frame; it is a useful
intermediate regulator, not an exactly Lorentz-invariant final theory.
\end{warningbox}

\Needspace{7\baselineskip}
\section{Schr\"odinger and Heisenberg descriptions}
In the Schr\"odinger picture, a state is a wavefunctional
\(\Psi[\varphi]=\langle\varphi|\Psi\rangle\), where \(\varphi(\bm x)\) labels a
field configuration, not a one-particle wavefunction. The regulated
operators act as
\[
 \phi(\bm x)\Psi=\varphi(\bm x)\Psi,\qquad
 \pi(\bm x)\Psi=-i\frac{\delta\Psi}{\delta\varphi(\bm x)}.
\]
The condition \(a(\bm p)\vac=0\) gives the Gaussian vacuum functional
\begin{equation}
 \Psi_0[\varphi]=\mathcal N\exp\left[-\frac12\int
 \frac{\dd^dp}{(2\pi)^d}E_{\bm p}\,
 \widetilde\varphi(\bm p)\widetilde\varphi(-\bm p)\right].
 \label{s-eq:gaussian}
\end{equation}
This is the product of oscillator ground states. Its normalization is
understood first with a finite regulator.

In the Heisenberg picture the state is fixed and
\begin{equation}
 O_H(t)=e^{iHt}O_S e^{-iHt},\qquad
 \partial_tO_H=i[H,O_H],\qquad
 a_H(\bm p,t)=e^{-iE_{\bm p}t}a(\bm p).
 \label{s-eq:Heisenberg}
\end{equation}
Thus \eqref{s-eq:realmode} already includes all free time evolution.
One may put the time dependence in \(a_H(t)\) or in the plane wave;
putting it in both would count it twice. The equal-time commutators hold
at every time by unitary conjugation. The interaction picture will later
separate the free evolution from interactions.

\Needspace{7\baselineskip}
\section{A complex scalar: two species and one charge}
\subsection{Independent canonical pairs}
Let \(\Phi=(\phi_1+i\phi_2)/\sqrt2\). Two canonically normalized real fields give
\begin{equation}
 \Lag=\partial_\mu\Phi^\dagger\partial^\mu\Phi-m^2\Phi^\dagger\Phi,
 \quad \Pi=\dot\Phi^\dagger,\quad \Pi^\dagger=\dot\Phi,
 \quad \mathcal H=\Pi^\dagger\Pi+\bm\nabla\Phi^\dagger\cdot\bm\nabla\Phi
                 +m^2\Phi^\dagger\Phi.
 \label{s-eq:complexaction}
\end{equation}
There is no extra factor \(1/2\) in this complex-field convention.
The two canonical brackets are
\([\Phi,\Pi]=i\delta^{(d)}\) and
\([\Phi^\dagger,\Pi^\dagger]=i\delta^{(d)}\); all others vanish at equal time.
Taking the adjoint of the first bracket also reverses its order, which
explains the same plus sign in the second.

The mode expansion and its adjoint are
\begin{align}
 \Phi(x)&=\int\dd\Pi_p[a(\bm p)e^{-ipx}+b^\dagger(\bm p)e^{ipx}],\notag\\
 \Phi^\dagger(x)&=\int\dd\Pi_p[b(\bm p)e^{-ipx}+a^\dagger(\bm p)e^{ipx}].
 \label{s-eq:complexmode}
\end{align}
Both \(a\) and \(b\) obey \eqref{s-eq:modeCCR}, and the two families commute.
Unlike a real scalar, Hermiticity no longer identifies the coefficients
in the two frequency sectors. After normal ordering,
\begin{equation}
 H_N=\int\dd\Pi_p E_{\bm p}(a^\dagger a+b^\dagger b),\qquad
 N=N_a+N_b=\int\dd\Pi_p(a^\dagger a+b^\dagger b).
 \label{s-eq:complexH}
\end{equation}
Both species have positive energy.

\subsection{Phase symmetry and antiparticles}
For \(\delta\Phi=-i\alpha\Phi\), Noether's theorem gives
\begin{equation}
 j^\mu=i(\Phi^\dagger\partial^\mu\Phi-
              (\partial^\mu\Phi^\dagger)\Phi),\qquad
 Q=\int\dd^dx:j^0:=\int\dd\Pi_p(a^\dagger a-b^\dagger b).
 \label{s-eq:charge}
\end{equation}
Indeed \([Q,a^\dagger]=a^\dagger\), \([Q,b^\dagger]=-b^\dagger\), and
\([Q,\Phi]=-\Phi\). The unitary \(e^{i\alpha Q}\) acts on the field as
\(\Phi\mapsto e^{-i\alpha}\Phi\). Reversing the phase parameter reverses
this naming convention, without changing the opposite charges.

Interactions respecting this \(U(1)\), such as \((\Phi^\dagger\Phi)^2\),
conserve \(Q\) while generally failing to conserve \(N\). Charge-neutral
particle--antiparticle pairs can be produced when energy and momentum
allow it. This does not imply that the vacuum of a stable isolated theory
spontaneously produces pairs: its energy cannot supply them.

\subsection{Charge conjugation}
With phases chosen consistently, the free theory admits a unitary
operator \(\Dc\) satisfying
\begin{equation}
 \Dc\Phi(x)\Dc^{-1}=\Phi^\dagger(x),\qquad
 \Dc a\Dc^{-1}=b,\quad \Dc b\Dc^{-1}=a,\qquad
 \Dc Q\Dc^{-1}=-Q,\quad \Dc H_N\Dc^{-1}=H_N.
 \label{s-eq:C}
\end{equation}
Charge conjugation exchanges the two species; it is distinct from a
continuous phase rotation. A definite nonzero-charge state is not also a
charge-conjugation eigenstate. In a theory with an electromagnetic
coupling, invariance requires \(A_\mu\mapsto-A_\mu\) along with
\(j^\mu\mapsto-j^\mu\); in a charge-conjugation-invariant vacuum this gives
a single photon \(C=-1\). This is a symmetry statement about the whole
theory, not a consequence of Lorentz covariance alone.



\Needspace{7\baselineskip}
\section{A fermionic oscillator and exclusion}
A single fermionic mode has the canonical anticommutation relations (CAR)
\begin{equation}
 \{c,c^\dagger\}=cc^\dagger+c^\dagger c=1,
 \qquad c^2=(c^\dagger)^2=0.
 \label{f-eq:oneCAR}
\end{equation}
Choose \(c\vac=0\), \(\langle0|0\rangle=1\), and define \(|1\rangle=c^\dagger\vac\).
The norm of \(|1\rangle\) is one and no second excitation is possible.
In this basis,
\[
 c=\begin{pmatrix}0&1\\0&0\end{pmatrix},\qquad
 c^\dagger=\begin{pmatrix}0&0\\1&0\end{pmatrix},\qquad
 n=c^\dagger c=\begin{pmatrix}0&0\\0&1\end{pmatrix}.
\]
Thus \(n^2=n\); its eigenvalues are zero and one. A symmetric choice of
energy origin is
\begin{equation}
 H=\frac\omega2(c^\dagger c-cc^\dagger)=\omega(n-\tfrac12),
 \qquad :H:=\omega n.
 \label{f-eq:fermionoscillator}
\end{equation}
The negative vacuum constant is not a negative excitation energy.
The energy difference is \(\omega>0\).

For distinct modes, all creation operators anticommute. Hence
\(c_r^\dagger c_s^\dagger\vac=-c_s^\dagger c_r^\dagger\vac\): the two-particle
wavefunction is antisymmetric. The many-mode space is
\(\mathcal F_F(\cH_1)=\bigoplus_n\bigwedge^n\cH_1\).
This exchange sign is different from the minus sign in a commutator.
For comparison, the bosonic relation is \(aa^\dagger-a^\dagger a=1\),
whereas \eqref{f-eq:oneCAR} has a plus sign.

\Needspace{7\baselineskip}
\section{Canonical quantization of the Dirac field}
\begin{keyidea}[Conventions inherited from the spinor lectures]
We use \(\eta=\diag(1,-1,-1,-1)\), \(p\cdot x=E_{\bm p}t-\bm p\cdot\bm x\),
and \(\dd\Pi_p=\dd^3p/[(2\pi)^3 2E_{\bm p}]\). The explicit matrices are
\[
 \gamma^0=\begin{pmatrix}0&1\\1&0\end{pmatrix},\qquad
 \gamma^i=\begin{pmatrix}0&-\sigma_i\\\sigma_i&0\end{pmatrix},\qquad
 \bar\Psi=\Psi^\dagger\gamma^0.
\]
Each displayed block is two dimensional. They satisfy
\(\{\gamma^\mu,\gamma^\nu\}=2\eta^{\mu\nu}\) and
\((\gamma^\mu)^\dagger=\gamma^0\gamma^\mu\gamma^0\).
The classical field is Grassmann odd; after quantization its coefficients
are operators on a fermionic Fock space. Ordinary mode spinors \(u_s,v_s\)
are complex-number columns, not Grassmann variables.
\end{keyidea}
\subsection{First-order dynamics and canonical momentum}
Start with the real-mass action
\begin{equation}
 S=\int\dd^4x\,\bar\Psi(i\gamma^\mu\partial_\mu-m)\Psi
   =\int\dd^4x\,[i\Psi^\dagger\dot\Psi-\Psi^\dagger h_D\Psi],
 \qquad h_D=-i\bm\alpha\cdot\bm\nabla+\beta m,
 \label{f-eq:action}
\end{equation}
where \(\alpha^i=\gamma^0\gamma^i\), \(\beta=\gamma^0\).
With the derivative convention corresponding to the displayed kinetic
term, the canonical momentum is \(\Pi_\Psi=i\Psi^\dagger\).
There is no independent momentum containing \(\dot\Psi^\dagger\).
These are constraints of a first-order action; one must not apply two
independent scalar canonical pairs. Equivalently, the reduced canonical
structure gives
\begin{equation}
 \{\Psi_\alpha(t,\bm x),\Psi_\beta^\dagger(t,\bm y)\}
 =\delta_{\alpha\beta}\delta^{(3)}(\bm x-\bm y),\qquad
 \{\Psi_\alpha,\Psi_\beta\}=\{\Psi_\alpha^\dagger,\Psi_\beta^\dagger\}=0.
 \label{f-eq:fieldCAR}
\end{equation}
The Hamiltonian is \(H=\int\dd^3x\,\Psi^\dagger h_D\Psi\).
The matrices \(\alpha^i,\beta\) are Hermitian, so this operator is formally
Hermitian with appropriate spatial boundary conditions.
The on-shell rewriting \(H=\int\dd^3x\,i\Psi^\dagger\dot\Psi\) follows from
the Dirac equation; it is not the starting definition of the Hamiltonian.

\subsection{Promoting the mode coefficients to operators}
Recall the spinors of Lecture 14:
\begin{gather}
 (\slashed p-m)u_s(p)=0,\qquad(\slashed p+m)v_s(p)=0,\notag\\
 u_s^\dagger u_t=v_s^\dagger v_t=2E_{\bm p}\delta_{st},\qquad
 \sum_su_s\bar u_s=\slashed p+m,\quad
 \sum_sv_s\bar v_s=\slashed p-m.
 \label{f-eq:spinorinput}
\end{gather}
Expand
\begin{equation}
 \Psi(x)=\sum_{s=1}^2\int\dd\Pi_p\,
 [b_s(\bm p)u_s(p)e^{-ipx}+d_s^\dagger(\bm p)v_s(p)e^{ipx}].
 \label{f-eq:modes}
\end{equation}
The two independent species have mode algebra
\begin{align}
 \{b_s(\bm p),b_t^\dagger(\bm q)\}
 &=\{d_s(\bm p),d_t^\dagger(\bm q)\}
  =(2\pi)^3 2E_{\bm p}\delta_{st}\delta^{(3)}(\bm p-\bm q),
 \label{f-eq:modeCAR}
\end{align}
with every other independent anticommutator zero, including mixed
particle--antiparticle ones. The vacuum obeys \(b_s\vac=d_s\vac=0\).
In a finite box divide the modes by \(\sqrt{2E_{\bm p}V}\) to obtain the
unit-normalized oscillators of \eqref{f-eq:oneCAR}.

\subsection{Checking the equal-time anticommutator}
It is useful to do this once rather than infer it by analogy with a scalar.
The mode algebra gives
\begin{align}
 \{\Psi_\alpha(t,\bm x),\Psi_\beta^\dagger(t,\bm y)\}
 &=\int\dd\Pi_p\sum_s\left[
 (u_su_s^\dagger)_{\alpha\beta}e^{i\bm p\cdot(\bm x-\bm y)}
 +(v_sv_s^\dagger)_{\alpha\beta}e^{-i\bm p\cdot(\bm x-\bm y)}\right].
 \label{f-eq:CARcheck}
\end{align}
Change \(\bm p\to-\bm p\) only in the second term, and set
\(p_-=(E_{\bm p},-\bm p)\). The numerator becomes
\begin{equation}
 (\slashed p+m)\gamma^0+(\slashed p_- -m)\gamma^0=2E_{\bm p}I_4.
 \label{f-eq:completeness}
\end{equation}
The mass and spatial-momentum terms cancel. The remaining Fourier
integral is exactly \eqref{f-eq:fieldCAR}. Omitting the reversed momentum
would spoil this check. The same-momentum statement
\(u_s^\dagger(\bm p)v_t(\bm p)=0\) is generally false; the needed identity is
\(u_s^\dagger(\bm p)v_t(-\bm p)=0\).

\Needspace{8\baselineskip}
\Needspace{7\baselineskip}
\section{Energy, charge, and particle states}
\subsection{Why antiparticles have positive excitation energy}
The positive-frequency term obeys
\(h_D[u_s(p)e^{i\bm p\cdot\bm x}]=E_{\bm p}u_s(p)e^{i\bm p\cdot\bm x}\).
For the negative-frequency term the eigenvalue is \(-E_{\bm p}\).
Spatial integration and the orthogonality just discussed remove the
mixed terms, giving
\begin{equation}
 H=\sum_s\int\dd\Pi_p E_{\bm p}
 [b_s^\dagger(\bm p)b_s(\bm p)-d_s(\bm p)d_s^\dagger(\bm p)].
 \label{f-eq:unordH}
\end{equation}
The minus sign here precedes use of the operator algebra. For a normalized
mode, \(-dd^\dagger=d^\dagger d-1\). Consequently
\begin{equation}
 H_N\equiv:H:=\sum_s\int\dd\Pi_p E_{\bm p}(b_s^\dagger b_s+d_s^\dagger d_s),
 \qquad \bm P=\sum_s\int\dd\Pi_p\bm p\,(b_s^\dagger b_s+d_s^\dagger d_s).
 \label{f-eq:normalH}
\end{equation}
In a box the un-subtracted vacuum energy is
\(-\sum_{\bm p,s}E_{\bm p}\): there are two spin labels and two charge
sectors, each with the oscillator offset \(-E/2\).
The normal-ordered excitation spectrum is nonnegative.
In particular \([H_N,b_s^\dagger]=E_{\bm p}b_s^\dagger\) and
\([H_N,d_s^\dagger]=E_{\bm p}d_s^\dagger\).

\begin{warningbox}[What goes wrong with a bosonic assignment?]
If both mode families were ordinary positive-norm bosonic oscillators,
\(-dd^\dagger=-d^\dagger d-1\); repeated \(d^\dagger\) excitations would
lower the energy without bound. Trying instead to impose the usual
local field commutator leads to a negative sign in the \(d\)-mode
commutator and hence negative-norm excitations for a vacuum annihilated
by \(d\). These are obstructions to combining positivity, locality,
and the Dirac transformation law. This free-field calculation motivates
the spin--statistics connection; it is not a proof of the full theorem.
\end{warningbox}

\subsection{A sum in the energy, a difference in the charge}
The continuous transformation \(\Psi\mapsto e^{-i\alpha}\Psi\) has current
\(j^\mu=\bar\Psi\gamma^\mu\Psi\). Before normal ordering, the charge
contains \(b^\dagger b+dd^\dagger\), because both \(u^\dagger u\) and
\(v^\dagger v\) are positive. Fermionic normal ordering then gives
\begin{equation}
 Q=\int\dd^3x:\Psi^\dagger\Psi:
   =\sum_s\int\dd\Pi_p(b_s^\dagger b_s-d_s^\dagger d_s).
 \label{f-eq:Q}
\end{equation}
Thus \([Q,b_s^\dagger]=b_s^\dagger\), \([Q,d_s^\dagger]=-d_s^\dagger\), and
\([Q,\Psi]=-\Psi\). As for the scalar field, \(e^{i\alpha Q}\) implements the
stated field phase. The normal-ordering rule includes a minus sign for
each exchange of odd operators; for example \(:dd^\dagger:=-d^\dagger d\).
Energy and charge distinguish the two sectors as follows:
\begin{center}
\begin{tabular}{lcc}
\toprule
State & Energy above vacuum & Charge\\
\midrule
\(b_s^\dagger(\bm p)\vac\) & \(E_{\bm p}\) & \(+1\)\\
\(d_s^\dagger(\bm p)\vac\) & \(E_{\bm p}\) & \(-1\)\\
\bottomrule
\end{tabular}
\end{center}
These generalized momentum states have norm
\((2\pi)^3 2E_{\bm p}\delta_{st}\delta^{(3)}(\bm p-\bm q)\).
Their matrix elements reproduce the classical wavefunctions:
\begin{equation}
 \vbra\Psi(x)b_s^\dagger(\bm p)\vac=u_s(p)e^{-ipx},\qquad
 \vbra\bar\Psi(x)d_s^\dagger(\bm p)\vac=\bar v_s(p)e^{-ipx}.
 \label{f-eq:matrixelements}
\end{equation}
There are two particle and two antiparticle polarizations at each
momentum, not four independent particle spin states.

\Needspace{7\baselineskip}
\section{The Dirac sea as a change of vacuum language}
With a finite number of regulated modes, write the coefficient of the
negative-frequency term as an annihilator \(f_s=d_s^\dagger\). Equation
\eqref{f-eq:unordH} then contains \(-E f_s^\dagger f_s\). If the reference
vacuum is empty in these modes, filling one lowers the energy. Since each
fermionic mode can be occupied only once, the energy is minimized by
filling all the negative-energy modes.

Relative to that filled state, a hole is created by removing an \(f\)
particle. In formulas,
\begin{equation}
 d_s=f_s^\dagger,\qquad d_s^\dagger=f_s,\qquad
 f_s^\dagger f_s=1-d_s^\dagger d_s.
 \label{f-eq:holes}
\end{equation}
A hole costs positive energy and carries the opposite charge. Subtracting
the energy and charge of the filled state gives precisely
\eqref{f-eq:normalH} and \eqref{f-eq:Q}. The particle--antiparticle vacuum
\(b_s\vac=d_s\vac=0\) encodes this directly.

The sea picture is a useful finite-regulator explanation of the signs,
not a literal finite-density medium in the Lorentz-invariant vacuum.
An infinite product of filled modes requires a limiting construction;
it need not be a vector in the original empty-sea Fock representation.
The physical lesson is that an operator algebra alone does not specify
a quantum theory's state or vacuum representation. Also, scalar
antiparticles exist without any exclusion principle or filled sea.

\section*{Summary}
\addcontentsline{toc}{section}{Summary}
Canonical quantization turns classical mode coefficients into operators on
Fock space. Covariant normalization fixes the factors of \(2E\).
Commutators build symmetric scalar states; anticommutators build
antisymmetric spinor states. Both theories have positive-energy
antiparticles. Energy adds the two charge sectors, while charge subtracts
them. Normal ordering removes a vacuum reference constant, and is not a
general replacement for renormalization. The finite-regulator hole picture
explains the fermionic signs without making a filled sea necessary for
scalar antiparticles. These operator constructions are the starting point
for the symmetry and locality tests of Lecture 16.

\newpage
\section*{Questions for review}
\addcontentsline{toc}{section}{Questions for review}
\begin{enumerate}
\item Derive the chain dispersion and identify the parameters held fixed
in its continuum limit.
\item Derive \eqref{s-eq:modeCCR} from the canonical brackets, then check the
inverse calculation directly in position space.
\item Show explicitly why the \(aa\) terms cancel in the free Hamiltonian.
\item Explain the difference between an occupation number in a box and a
continuum momentum density.
\item Derive the invariant one-particle measure and the wave-packet norm.
Why is \(\phi(x)\vac\) not a normalized localized particle?
\item Separate the infinite-volume factor in \(E_0\) from the ultraviolet
divergence of \(\rho_0\). Which changes of the energy zero affect free evolution?
\item Verify the Gaussian vacuum functional by applying \(a(\bm p)\).
\item Derive the complex-scalar Hamiltonian and charge. Why does normal
ordering give a sum in one and a difference in the other?
\item Compare particle-number conservation, charge conservation, and
charge-conjugation symmetry.
\item What must be specified before a claim of cutoff independence is meaningful?
\item Prove \(n^2=n\) for one fermionic oscillator and derive its two energies.
\item Why is \(\Pi_\Psi=i\Psi^\dagger\), rather than \(\dot\Psi^\dagger\)?
What feature of the action makes the canonical system constrained?
\item Verify the equal-time CAR using the two spin sums and the change
\(\bm p\to-\bm p\).
\item Derive the un-normal-ordered Hamiltonian, then track the reordering
sign in the antiparticle sector.
\item Repeat the argument for \(Q\), explaining why its intermediate sign
differs from the Hamiltonian's.
\item Describe the hole transformation for a finite set of modes.
What extra issue arises for an infinite sea?
\end{enumerate}

\section*{Sources and reading}
\addcontentsline{toc}{section}{Sources and reading}
Based on the supplied scalar-field slides and the quantization/Dirac-sea
sections of the supplied Dirac-field slides. Their discrete-symmetry and
locality discussions continue in Lecture 16. For comparison, see D.~Tong,
\emph{Quantum Field Theory}, Chapters 2 and 5,
\url{https://www.damtp.cam.ac.uk/user/tong/qft.htm}.
\end{document}
